\documentclass[11pt]{article}

\usepackage[a4paper,margin=1in]{geometry}
\usepackage{amsmath,amssymb}
\usepackage{hyperref}

\hypersetup{
    colorlinks=true,
    linkcolor=blue,
    citecolor=blue,
    urlcolor=blue
}

\title{Pushforward Measures and Change of Variables}
\author{}
\date{}

\begin{document}

\maketitle

\section*{Pushforward}

Let
\[
F:X\to Y
\]
be a measurable map and let $\mu$ be a probability measure on $X$. The
pushforward measure
\[
F_\#\mu
\]
is the probability measure on $Y$ defined by
\[
(F_\#\mu)(B)=\mu(F^{-1}(B))
\]
for every measurable set $B\subseteq Y$.

\section*{Defining Integral Property}

The defining integral property of the pushforward is:
\[
\int_Y \Phi(y)\,d(F_\#\mu)(y)
=
\int_X \Phi(F(x))\,d\mu(x)
\]
for every integrable function $\Phi:Y\to\mathbb R$.

\section*{Why This Is a Change of Variables}

The left-hand side integrates over the pushed-forward variable $y$. The
right-hand side pulls the same function back to the original variable $x$ by
composing with $F$.

Thus the notation
\[
y=F(x)
\]
should be understood measure-theoretically as
\[
d(F_\#\mu)(y)
\quad\text{is the distribution of}\quad
F(x)\ \text{when}\ x\sim\mu.
\]

\section*{Application to the Geodesic Proof}

In the Wasserstein geodesic proof, we define
\[
F_{s,t}(x)=(T_s(x),T_t(x))
\]
and
\[
\gamma_{s,t}=(F_{s,t})_\#\mu.
\]
Taking
\[
\Phi(a,b)=|a-b|^2
\]
gives
\[
\int_{\mathbb R^d\times\mathbb R^d}
|a-b|^2\,d\gamma_{s,t}(a,b)
=
\int_{\mathbb R^d}
|T_s(x)-T_t(x)|^2\,d\mu(x).
\]
The apparent double integral becomes a single integral because
$\gamma_{s,t}$ is not arbitrary: it is generated by one original particle
label $x$.

\end{document}
